\documentclass[12pt, pdf, xcolor={table, dvipsnames}, paperheight=8cm,paperwidth=12cm]{beamer}
\definecolor{tugreen}{RGB}{132,184,25}
\usepackage{../assets/themes/beamerthemeTUDortmund}
\usepackage{../assets/themes/beamercolorthemetudortmund}

%\expandafter\def\csname ver@transparent.sty\endcsname{}
%\usepackage[edges]{forest}
%\usepackage{pgf-umlsd}
%\usepackage{dsfont}
\usepackage[utf8]{inputenc}
\usepackage[ngerman]{babel}
\usepackage[style=german,babel]{csquotes}
\usepackage{glossaries}
\makeindex
%\usepackage{amsmath}
%\usepackage{graphicx}
%\usepackage{pgfplots} 
%\usepackage{fontawesome}
%\usepackage{import}
%\usepackage{listings}
%\usepackage{titleref}
%\usepackage{caption}
%\usepackage{tikz}
%\usepackage{../assets/tikz-uml}
%\usepackage{biblatex}
%\addbibresource{literatur.bib}
%\usetikzlibrary{backgrounds,shapes.multipart}
%\pgfdeclarelayer{myback}
%\pgfsetlayers{background,main,myback}
%\makeatletter
%\newcommand*{\currentname}{\TR@currentTitle}
%\newcommand*{\overlaynumber}{\number\beamer@slideinframe}
%\makeatother
%\usetikzlibrary{positioning,arrows,automata,shapes.multipart,calc, fit, decorations.shapes,shadows}
%\tikzset{
%  mylabel/.style = {font=\footnotesize, midway, fill=white, anchor=center},
%  decorate sep/.style 2 args=
%    {decorate,decoration={shape backgrounds,shape=circle,shape size=#1,shape sep=#2}},
%  invisible/.style={opacity=0, text opacity=0},
%  visible on/.style={alt={#1{}{invisible}}},
%  alt/.code args={<#1>#2#3}{%
%    \alt<#1>{\pgfkeysalso{#2}}{\pgfkeysalso{#3}} % \pgfkeysalso doesn't change the path
%  },
%  onslide/.code args={<#1>#2}{%
%  \only<#1>{\pgfkeysalso{#2}} % \pgfkeysalso doesn't change the path
%    }
%}
%\forestset{
%  step/.style={
%    /tikz/visible on={<+->},
%    edge={/tikz/visible on={<+->}},
%    for descendants={
%        /tikz/visible on={<.->},
%        edge={/tikz/visible on={<.->}}
%    }
%  },
%  highlight/.style={
%    for tree={
%      /tikz/onslide=<.>{draw=red},
%      edge+={/tikz/onslide=<.>{draw=red}}
%    }
%  }
%}
%
\author{Dennis Ochocki}
\title[]{Sensordatenverarbeitung durch metamodelierte best-of-breed Lösung}
%\subtitle{Grafana, Prometheus und Docker}
%\titlegraphic{
%\includegraphics[width=0.35\textwidth]{../assets/themes/img/tud_logo_rgb}\hspace*{3.75cm}~%
%   \includegraphics[width=0.25\textwidth]{../assets/themes/img/fakultaet.pdf}
%}
%\logo{}
%\institute{}
%\date{\today}
%\subject{}
%\setbeamertemplate{navigation symbols}{}
%\setbeamertemplate{frametitle}{%
%    \nointerlineskip%
%    \begin{beamercolorbox}[wd=\paperwidth,ht=2.6ex,dp=1.6ex]{frametitle}
%        \hspace*{1ex}\insertframetitle%
%    \end{beamercolorbox}%
%}
%\setlength{\skip\footins}{0cm}
%\addtobeamertemplate{frametitle}{}{%
%\begin{tikzpicture}[remember picture,overlay]
%\node[anchor=north east,yshift=0pt] at (current page.north east) {\includesvg[height=0.7cm]{../themes/img/fakultaet}};
%\end{tikzpicture}}
%
\AtBeginSection[]{
  \begin{frame}
  \vfill
  \centering
  \begin{beamercolorbox}[sep=8pt,center,shadow=false,rounded=false]{title}
    \usebeamerfont{title}\insertsectionhead\par%
  \end{beamercolorbox}
  \vfill
  \end{frame}
}

\import{abbr}

\begin{document}
	
\maketitle

\begin{frame}{Übersicht}
	\tableofcontents
\end{frame}

\section{Motivation}
\begin{frame}{Motivation}
	Smart Home ist immer präsenter und dank \gls{iot} erst möglich. Immer mehr Geräte, sog. \enquote{Smart Devices}, kommen auf den Markt, die in solche Systeme integriert werden können. 
	
	\begin{block}{Was sind die Schwierigkeiten?}
		\begin{itemize}
			\item Bei vielen Smart Devices ist es schwer den Überblick zu behalten
			\item Assoziationen von Geräte-IP zur physikalischer Position muss gemerkt werden
			\item Eigenschaften des Sensors müssen manuell notiert werden
			\item
		\end{itemize}
	\end{block}
\end{frame}


\section{Grundlagen}
\begin{frame}{Grundlagen}
	\begin{block}{Was ist Smart Home:}
		\begin{itemize}
			\item Viele verschiedene Geräte in einem Netzwerk
			\item Jedes Gerät ist Beobachter und/oder Akteur
			\item Dient der Qualitätssteigerung der vernetzten Örtlichkeit
		\end{itemize}
	\end{block}
\end{frame}

\begin{frame}{Senoren}
\end{frame}

\begin{frame}{Prometheus}
	Prometheus is an open-source systems monitoring and alerting toolkit
\end{frame}

\begin{frame}{Grafana}
\end{frame}

\begin{frame}{Cinco}
\end{frame}

\section{Technologien}
\begin{frame}{CNNM}
\end{frame}


% \section{Motivation}
% \begin{frame}{\secname}
% \begin{exampleblock}{Konzept}
% \onslide<+->
% \begin{itemize}[opacity=.2]
% \item<+-> Haben reelle Situation aus dem Straßenverkehr
% \item<+-> Wollen diese so einfach wie möglich textuel beschreiben
% \item<+-> Gängige Formate wie z.B. OpenScenario sehr komplex
% \item<+-> Hoher Grad von Abstraktion gewünscht
% \item<+-> Geht das überhaupt ohne essentiellen Informationsverlust?
% \end{itemize}
% \end{exampleblock}
% \end{frame}

\begin{frame}[plain,allowframebreaks]{Verwandte Arbeiten}
	%\nocite{*}
	%\printbibliography[heading=none, keyword={related}]
\end{frame}

\begin{frame}[plain,allowframebreaks]{Mögliche Literatur}
	%\printbibliography[heading=none, keyword={possible}]
\end{frame}
\end{document}